\documentclass[10pt,a4paper]{amsart}
\usepackage{amsmath,amssymb,amsthm,graphicx}
% Retained minimal preamble for the Stacks Project chapter flat.tex (More on Flatness).
% Verbatim excerpt of the chapter preamble: Xy-pic setup (native preamble lines 14-18) and the native
% theorem declarations with the mathematical macros (native preamble lines 154-217). Copied mechanically;
% nothing added, removed or redefined. The wrapper class and the amsmath/amssymb/amsthm/graphicx base
% packages are supplied by the assembler, so the class/externaldocument part of the native preamble is omitted.
% For commutative diagrams we use Xy-pic
\usepackage[all]{xy}

% We use 2cell for 2-commutative diagrams.
\xyoption{2cell}
% Theorem environments.
%
\theoremstyle{plain}
\newtheorem{theorem}[subsection]{Theorem}
\newtheorem{proposition}[subsection]{Proposition}
\newtheorem{lemma}[subsection]{Lemma}

\theoremstyle{definition}
\newtheorem{definition}[subsection]{Definition}
\newtheorem{example}[subsection]{Example}
\newtheorem{exercise}[subsection]{Exercise}
\newtheorem{situation}[subsection]{Situation}

\theoremstyle{remark}
\newtheorem{remark}[subsection]{Remark}
\newtheorem{remarks}[subsection]{Remarks}

\numberwithin{equation}{subsection}

% Macros
%
\def\lim{\mathop{\mathrm{lim}}\nolimits}
\def\colim{\mathop{\mathrm{colim}}\nolimits}
\def\Spec{\mathop{\mathrm{Spec}}}
\def\Hom{\mathop{\mathrm{Hom}}\nolimits}
\def\Ext{\mathop{\mathrm{Ext}}\nolimits}
\def\SheafHom{\mathop{\mathcal{H}\!\mathit{om}}\nolimits}
\def\SheafExt{\mathop{\mathcal{E}\!\mathit{xt}}\nolimits}
\def\Sch{\mathit{Sch}}
\def\Mor{\mathop{\mathrm{Mor}}\nolimits}
\def\Ob{\mathop{\mathrm{Ob}}\nolimits}
\def\Sh{\mathop{\mathit{Sh}}\nolimits}
\def\NL{\mathop{N\!L}\nolimits}
\def\CH{\mathop{\mathrm{CH}}\nolimits}
\def\proetale{{pro\text{-}\acute{e}tale}}
\def\etale{{\acute{e}tale}}
\def\QCoh{\mathit{QCoh}}
\def\Ker{\mathop{\mathrm{Ker}}}
\def\Im{\mathop{\mathrm{Im}}}
\def\Coker{\mathop{\mathrm{Coker}}}
\def\Coim{\mathop{\mathrm{Coim}}}

% Boxtimes
%
\DeclareMathSymbol{\boxtimes}{\mathbin}{AMSa}{"02}

%
% Macros for moduli stacks/spaces
%
\def\QCohstack{\mathcal{QC}\!\mathit{oh}}
\def\Cohstack{\mathcal{C}\!\mathit{oh}}
\def\Spacesstack{\mathcal{S}\!\mathit{paces}}
\def\Quotfunctor{\mathrm{Quot}}
\def\Hilbfunctor{\mathrm{Hilb}}
\def\Curvesstack{\mathcal{C}\!\mathit{urves}}
\def\Polarizedstack{\mathcal{P}\!\mathit{olarized}}
\def\Complexesstack{\mathcal{C}\!\mathit{omplexes}}
% \Pic is the operator that assigns to X its picard group, usage \Pic(X)
% \Picardstack_{X/B} denotes the Picard stack of X over B
% \Picardfunctor_{X/B} denotes the Picard functor of X over B
\def\Pic{\mathop{\mathrm{Pic}}\nolimits}
\def\Picardstack{\mathcal{P}\!\mathit{ic}}
\def\Picardfunctor{\mathrm{Pic}}
\def\Deformationcategory{\mathcal{D}\!\mathit{ef}}

% --- appended 2026-10-12 (Chief): the project own definitions of the reference/slogan/history
% display environments, copied mechanically from the shared Stacks preamble lines 8-12
% (\usepackage{verbatim} plus the three \newenvironment lines, which the chapter preamble keeps
% verbatim from the project preamble). Nothing added, removed or redefined: these environments
% carry no mathematical content and the project itself typesets them as comments. Without these
% three lines a counted statement whose own text uses them fails with Environment slogan undefined.
\usepackage{verbatim}
\newenvironment{reference}{\comment}{\endcomment}
\newenvironment{slogan}{\comment}{\endcomment}
\newenvironment{history}{\comment}{\endcomment}
\usepackage{cleveref}
\crefname{originalreference}{original reference}{original references}
\Crefname{originalreference}{Original reference}{Original references}
\makeatletter
\newlabel{situation-existence}{{0CTC}{source}{}{sourcepointer-1}{}}
\newlabel{situation-existence@cref}{{[originalreference][0][]0CTC}{[1][0][]source}{}{}{}}
\makeatother
\crefname{originallemma}{original lemma}{original lemmas}
\Crefname{originallemma}{Original lemma}{Original lemmas}
\makeatletter
\newlabel{lemma-compute-what-it-should-be}{{0CTD}{source}{}{sourcepointer-2}{}}
\newlabel{lemma-compute-what-it-should-be@cref}{{[originallemma][0][]0CTD}{[1][0][]source}{}{}{}}
\makeatother
\crefname{originallemma}{original lemma}{original lemmas}
\Crefname{originallemma}{Original lemma}{Original lemmas}
\makeatletter
\newlabel{lemma-finitely-presented}{{0CTH}{source}{}{sourcepointer-3}{}}
\newlabel{lemma-finitely-presented@cref}{{[originallemma][0][]0CTH}{[1][0][]source}{}{}{}}
\makeatother
\crefname{originallemma}{original lemma}{original lemmas}
\Crefname{originallemma}{Original lemma}{Original lemmas}
\makeatletter
\newlabel{more-algebra-lemma-limit-henselian}{{0CT7}{source}{}{sourcepointer-4}{}}
\newlabel{more-algebra-lemma-limit-henselian@cref}{{[originallemma][0][]0CT7}{[1][0][]source}{}{}{}}
\makeatother
\crefname{originallemma}{original lemma}{original lemmas}
\Crefname{originallemma}{Original lemma}{Original lemmas}
\makeatletter
\newlabel{more-morphisms-lemma-split-off-proper-part-henselian}{{0CT9}{source}{}{sourcepointer-5}{}}
\newlabel{more-morphisms-lemma-split-off-proper-part-henselian@cref}{{[originallemma][0][]0CT9}{[1][0][]source}{}{}{}}
\makeatother
\title{After shrinking the open the support of the constructed module is proper over the base}
\author{The Stacks Project Authors}
\date{stacks-project:flat.tex:local-snapshot-20261011; GNU FDL 1.2 or later}
\begin{document}
\maketitle
\noindent\textit{Editorial note.} All delivered mathematical text is the author's own, copied from the pinned Stacks Project chapter file flat.tex; no statement or proof step was rewritten, reordered or removed.\par\medskip
\noindent\textit{Editorial note.} The preamble is the chapter's own retained minimal preamble (sources/stacks-flat/retained-preamble.tex), a verbatim excerpt of the native Xy-pic setup and of the native theorem declarations and mathematical macros, followed by the project own definitions of the reference, slogan and history display environments (verbatim package plus three newenvironment lines) that the retained file carries at its end and that the project itself typesets as comments. No macro was invented or redefined, and the wrapper class and the amsmath/amssymb/amsthm/graphicx base packages are supplied by the assembler.\par\medskip
\noindent\textit{Editorial note.} The counted statement and proof contain no cross reference, no citation and no included file, so the delivered body is self contained.\par\medskip
\noindent\textit{Editorial note.} Canonical location of the counted proposition: Stacks tag 0CTI (sources/stacks-flat/pinned-original-source.tex line 8245).\par\medskip
\noindent\textit{Allowed standard prerequisites (editorial).}\begin{itemize}
\item The pair (A, Ker(A -> A\_n)) is henselian. More on Algebra, Lemma labelled more-algebra-lemma-limit-henselian, cited at proof lines 8255-8258. Makes the henselian splitting lemma applicable.
\item The finite presentation of F and the identification of the reductions of its support. Same chapter, Lemma labelled lemma-finitely-presented part (3), cited at proof lines 8260-8264. Identifies Z x Spec(A\_n) with the support of F\_n and shows Z is closed.
\item Splitting off the proper part over a henselian pair. More on Morphisms, Lemma labelled more-morphisms-lemma-split-off-proper-part-henselian, cited at proof lines 8266-8270. Produces Z\_1 proper over A with the required reductions, leaving Z\_2 off the closed fibre.
\end{itemize}\medskip
\ifdefined\hypertarget\hypertarget{sourcepointer-1}{}\fi
\noindent\textit{Source reference.} Original reference 0CTC: sources/stacks-flat/pinned-original-source.tex, source lines 8036--8036. Same-chapter pointer into the author environment situation at source line 8036; invoked at proof line None of the counted proof. It is retained only as a navigation pointer; every load-bearing step is inside the delivered span.\par\medskip
\ifdefined\hypertarget\hypertarget{sourcepointer-2}{}\fi
\noindent\textit{Source reference.} Original lemma 0CTD: sources/stacks-flat/pinned-original-source.tex, source lines 8062--8062. Same-chapter pointer into the author environment lemma at source line 8062; invoked at proof line None of the counted proof. It is retained only as a navigation pointer; every load-bearing step is inside the delivered span.\par\medskip
\ifdefined\hypertarget\hypertarget{sourcepointer-3}{}\fi
\noindent\textit{Source reference.} Original lemma 0CTH: sources/stacks-flat/pinned-original-source.tex, source lines 8181--8181. Same-chapter pointer into the author environment lemma at source line 8181; invoked at proof line 8259 of the counted proof. It is retained only as a navigation pointer; every load-bearing step is inside the delivered span.\par\medskip
\ifdefined\hypertarget\hypertarget{sourcepointer-4}{}\fi
\noindent\textit{Source reference.} Original lemma 0CT7: sources/stacks-chow/tags.txt, source lines 16057--16057. Cross-chapter standard prerequisite located in the project tag list staged in the local snapshot; invoked at proof line 8255 of the counted proof. It is retained only as a navigation pointer; every load-bearing step is inside the delivered span.\par\medskip
\ifdefined\hypertarget\hypertarget{sourcepointer-5}{}\fi
\noindent\textit{Source reference.} Original lemma 0CT9: sources/stacks-chow/tags.txt, source lines 16059--16059. Cross-chapter standard prerequisite located in the project tag list staged in the local snapshot; invoked at proof line 8264 of the counted proof. It is retained only as a navigation pointer; every load-bearing step is inside the delivered span.\par\medskip
\begin{lemma}
\label{lemma-proper-support}
In Situation \ref{situation-existence} let $K$ be as in
Lemma \ref{lemma-compute-what-it-should-be}. Let $W \subset X$
be as in Lemma \ref{lemma-finitely-presented}.
Set $\mathcal{F} = H^0(K)|_W$. Then, after possibly shrinking the open $W$,
the support of $\mathcal{F}$ is proper over $A$.
\end{lemma}

\begin{proof}
Fix $n \geq 1$. Let $I_n = \Ker(A \to A_n)$.
By More on Algebra, Lemma \ref{more-algebra-lemma-limit-henselian}
the pair $(A, I_n)$ is henselian.
Let $Z \subset W$ be the support of $\mathcal{F}$.
This is a closed subset as $\mathcal{F}$ is of finite presentation.
By part (3) of Lemma \ref{lemma-finitely-presented}
we see that $Z \times_{\Spec(A)} \Spec(A_n)$
is equal to the support of $\mathcal{F}_n$ and hence
proper over $\Spec(A/I)$.
By More on Morphisms, Lemma
\ref{more-morphisms-lemma-split-off-proper-part-henselian}
we can write $Z = Z_1 \amalg Z_2$ with $Z_1, Z_2$ open and
closed in $Z$, with $Z_1$ proper
over $A$, and with $Z_1 \times_{\Spec(A)} \Spec(A/I_n)$
equal to the support of $\mathcal{F}_n$.
In other words, $Z_2$ does not meet $X_0$.
Hence after replacing $W$ by $W \setminus Z_2$ we obtain the lemma.
\end{proof}

\end{document}
